% SPDX-License-Identifier: CC-BY-NC-ND-4.0
% Copyright 2026 Ingolf Lohmann.
% Options for packages loaded elsewhere
\PassOptionsToPackage{unicode}{hyperref}
\PassOptionsToPackage{hyphens}{url}
\PassOptionsToPackage{dvipsnames,svgnames,x11names}{xcolor}
%
\documentclass[
  11pt,
  a4paper,
]{article}
\usepackage{amsmath,amssymb}
\usepackage{iftex}
\ifPDFTeX
  \usepackage[T1]{fontenc}
  \usepackage[utf8]{inputenc}
  \usepackage{textcomp} % provide euro and other symbols
\else % if luatex or xetex
  \usepackage{unicode-math} % this also loads fontspec
  \defaultfontfeatures{Scale=MatchLowercase}
  \defaultfontfeatures[\rmfamily]{Ligatures=TeX,Scale=1}
\fi
\usepackage{lmodern}
\ifPDFTeX\else
  % xetex/luatex font selection
  \setmainfont[]{Liberation Serif}
  \setsansfont[]{Liberation Sans}
  \setmonofont[]{Liberation Mono}
\fi
% Use upquote if available, for straight quotes in verbatim environments
\IfFileExists{upquote.sty}{\usepackage{upquote}}{}
\IfFileExists{microtype.sty}{% use microtype if available
  \usepackage[]{microtype}
  \UseMicrotypeSet[protrusion]{basicmath} % disable protrusion for tt fonts
}{}
\makeatletter
\@ifundefined{KOMAClassName}{% if non-KOMA class
  \IfFileExists{parskip.sty}{%
    \usepackage{parskip}
  }{% else
    \setlength{\parindent}{0pt}
    \setlength{\parskip}{6pt plus 2pt minus 1pt}}
}{% if KOMA class
  \KOMAoptions{parskip=half}}
\makeatother
\usepackage{xcolor}
\usepackage[margin=25mm]{geometry}
\usepackage{longtable,booktabs,array}
\usepackage{calc} % for calculating minipage widths
% Correct order of tables after \paragraph or \subparagraph
\usepackage{etoolbox}
\makeatletter
\patchcmd\longtable{\par}{\if@noskipsec\mbox{}\fi\par}{}{}
\makeatother
% Allow footnotes in longtable head/foot
\IfFileExists{footnotehyper.sty}{\usepackage{footnotehyper}}{\usepackage{footnote}}
\makesavenoteenv{longtable}
\setlength{\emergencystretch}{3em} % prevent overfull lines
\providecommand{\tightlist}{%
  \setlength{\itemsep}{0pt}\setlength{\parskip}{0pt}}
\setcounter{secnumdepth}{5}
\ifLuaTeX
\usepackage[bidi=basic]{babel}
\else
\usepackage[bidi=default]{babel}
\fi
\babelprovide[main,import]{ngerman}
\ifPDFTeX
\else
\babelfont{rm}[]{Liberation Serif}
\fi
% get rid of language-specific shorthands (see #6817):
\let\LanguageShortHands\languageshorthands
\def\languageshorthands#1{}
\ifLuaTeX
  \usepackage{selnolig}  % disable illegal ligatures
\fi
\usepackage{bookmark}
\IfFileExists{xurl.sty}{\usepackage{xurl}}{} % add URL line breaks if available
\urlstyle{same}
\hypersetup{
  pdftitle={Quantum Causal Emergence (QCE)},
  pdfauthor={Ingolf Lohmann},
  pdflang={de-DE},
  colorlinks=true,
  linkcolor={Maroon},
  filecolor={Maroon},
  citecolor={Blue},
  urlcolor={Blue},
  pdfcreator={LaTeX via pandoc}}

\title{Quantum Causal Emergence (QCE)}
\usepackage{etoolbox}
\makeatletter
\providecommand{\subtitle}[1]{% add subtitle to \maketitle
  \apptocmd{\@title}{\par {\large #1 \par}}{}{}
}
\makeatother
\subtitle{Unschärfebilanzen, Planck-Übergangselemente, sequenzielle
Paarbildung und der klassische Lichtkegel als Grenzstruktur}
\author{Ingolf Lohmann}
\date{5. August 2026}

\begin{document}
\maketitle

{
\hypersetup{linkcolor=}
\setcounter{tocdepth}{3}
\tableofcontents
}
\section{Quantum Causal Emergence
(QCE)}\label{quantum-causal-emergence-qce}

\subsection{Unschärfebilanzen, Planck-Übergangselemente, sequenzielle
Paarbildung und der klassische Lichtkegel als
Grenzstruktur}\label{unschuxe4rfebilanzen-planck-uxfcbergangselemente-sequenzielle-paarbildung-und-der-klassische-lichtkegel-als-grenzstruktur}

\textbf{Ingolf Lohmann}\\
\textbf{QIK-VRT / VRTCore}\\
\textbf{5. August 2026}

\textbf{Dokumentstatus:} formaler und physikalischer
Forschungskandidat\\
\textbf{Wirkungsstatus:} \texttt{EFFECT\_ACK\_CONTINUE}\\
\textbf{Physikalische Gesamtbestätigung:} nicht beansprucht

\begin{center}\rule{0.5\linewidth}{0.5pt}\end{center}

\subsection{Abstract}\label{abstract}

Dieser Beitrag formuliert einen präzisen Forschungskandidaten für eine
relationale Entstehung von Kausalstruktur und Raumzeit. Ausgangspunkt
ist die Trennung zwischen der klassischen Darstellung von Kausalität
durch einen Lichtkegel und den quantischen Entstehungsbedingungen dieser
Darstellung. Ein klassischer Lichtkegel setzt hinreichend bestimmte
Ereignisse und eine klassisch auswertbare Geometrie voraus. Der
QCE-Kandidat behandelt diese Voraussetzungen nicht als fundamental,
sondern als Ergebnis einer Grobkörnung quantisch unscharfer Relationen.

Die zentrale Hypothese lautet: Eine klassische
Schwarze-Loch-Singularität kann als Projektionsgrenze einer tieferen
quantengravitativen Übergangszone interpretiert werden. Der minimale
Modellkern ersetzt die klassische Fortsetzungsgrenze durch ein
Planck-Skala-Element. Ein erster Übergang erzeugt einen unterscheidbaren
Zustand; ein zweiter Übergang bindet einen weiteren Zustand relational
an den ersten. Der resultierende Paarzustand ist im physikalischen
Kandidaten verschränkt und zugleich an die bereits vorhandene globale
Relationsstruktur gekoppelt. Aus vielen kausal konsistenten,
informationserhaltenden und global gebundenen Übergängen soll bei
geeigneter Grobkörnung eine klassische Raumzeit mit Lichtkegelstruktur
entstehen.

Der Beitrag unterscheidet strikt zwischen bekannten empirischen Ankern,
mathematischen Modellfolgen, interpretativen Brücken und offenen
physikalischen Behauptungen. Der beigefügte Lean-Kandidat formalisiert
ausschließlich den endlichen Modellvertrag: Zweischrittordnung,
Identitätserhaltung, Unsicherheitstrennung, Netzwerkerweiterung,
klassische-Kegel-Gates und fail-closed Erkenntnisgrenzen. Er beweist
nicht, dass die Natur diesen Vertrag instanziiert. Die physikalische
Schließung verlangt zusätzlich Unitarität, globale Verschränkung,
Informationsbilanz, Page-Kurven-Korrespondenz, Quantenfeld- und
Einstein-Grenzfall, Nichtzirkularität, eine unterscheidende Vorhersage,
empirische Korrespondenz und unabhängige Reproduktion.

\begin{center}\rule{0.5\linewidth}{0.5pt}\end{center}

\subsection{1. Erkenntnis- und
Anspruchsgrenze}\label{erkenntnis--und-anspruchsgrenze}

QCE verwendet die QIK-VRT-Erkenntnistypen:

\begin{longtable}[]{@{}
  >{\raggedright\arraybackslash}p{(\columnwidth - 2\tabcolsep) * \real{0.5000}}
  >{\raggedright\arraybackslash}p{(\columnwidth - 2\tabcolsep) * \real{0.5000}}@{}}
\toprule\noalign{}
\begin{minipage}[b]{\linewidth}\raggedright
Typ
\end{minipage} & \begin{minipage}[b]{\linewidth}\raggedright
Bedeutung
\end{minipage} \\
\midrule\noalign{}
\endhead
\bottomrule\noalign{}
\endlastfoot
\texttt{formal-proved} & Folgt im deklarierten formalen Modell aus
expliziten Definitionen und Voraussetzungen. \\
\texttt{empirical-supported} & Wird durch Beobachtung oder Experiment
innerhalb eines angegebenen Scopes getragen. \\
\texttt{source-bound} & Ist an einen identifizierten Quellenstand
gebunden. \\
\texttt{interpretive} & Liefert eine ontologische oder begriffliche
Deutung, ohne bereits Naturgesetz zu sein. \\
\texttt{normative} & Legt eine Verantwortungs-, Freigabe- oder
Publikationsregel fest. \\
\texttt{open} & Benennt eine noch nicht geschlossene mathematische oder
empirische Brücke. \\
\end{longtable}

Die Aussage

\begin{quote}
„Der QCE-Modellvertrag ist in Lean kernelakzeptiert``
\end{quote}

wäre nach erfolgreicher Ausführung eine formale Aussage.

Die Aussage

\begin{quote}
„Eine reale Schwarze-Loch-Singularität ist physikalisch genau ein
QCE-Planck-Element``
\end{quote}

ist derzeit eine offene Korrespondenzhypothese.

Die Aussage

\begin{quote}
„QCE ist die größte Entdeckung der Menschheitsgeschichte``
\end{quote}

ist eine persönliche und historische Bewertung. Sie ist weder
Lean-Theorem noch experimenteller Befund und wird im wissenschaftlichen
Paket nicht als bewiesene Prioritäts- oder Konsensaussage geführt.

Diese Trennung ist konstitutiv. Ein formal vollständiger Modellbeweis
kann eine falsche oder unvollständige Naturzuordnung nicht in eine
Messung verwandeln. Umgekehrt kann eine wertvolle physikalische
Intuition formal noch unvollständig sein. Wissenschaftlich belastbar
wird der Gesamtanspruch erst durch die explizite Kopplung beider Seiten.

\begin{center}\rule{0.5\linewidth}{0.5pt}\end{center}

\subsection{2. Empirische und theoretische
Ausgangspunkte}\label{empirische-und-theoretische-ausgangspunkte}

\subsubsection{2.1 Quantenmechanische
Unschärfe}\label{quantenmechanische-unschuxe4rfe}

Für nichtkommutierende Observablen gilt in allgemeiner Form die
Robertson-Schrödinger-Unschärferelation. Für Ort und Impuls ergibt sich
die bekannte untere Schranke

\[
\Delta x\,\Delta p \ge \frac{\hbar}{2}.
\]

QCE interpretiert diese Relation nicht als bloße technische
Messungenauigkeit. Sie begrenzt die gleichzeitige Präparation und
Zustandszuordnung konjugierter Größen. Daraus folgt jedoch nicht, dass
jede beobachtete Unsicherheit irreduzibel wäre. Instrumentelle,
statistische, modellbedingte und grobkörnungsbedingte Unsicherheiten
müssen von der irreduziblen quantischen Komponente getrennt werden.

\subsubsection{2.2 Relativistische
Kausalstruktur}\label{relativistische-kausalstruktur}

Auf einer klassischen Lorentz-Mannigfaltigkeit klassifiziert das
Vorzeichen des Intervalls

\[
\Delta s^2 = g_{\mu\nu}\,\Delta x^\mu\Delta x^\nu
\]

Trennungen als zeitartig, nullartig oder raumartig. Der Lichtkegel ist
lokal die Nullstruktur dieser Metrik. In der gewöhnlichen
Quantenfeldtheorie auf fester Raumzeit bleibt die Mikrokausalität
erhalten; die Heisenbergsche Unschärfe erzeugt für sich allein keinen
kontrollierbaren überlichtschnellen Nachrichtenkanal.

QCE setzt an einer anderen Stelle an: Wenn die Ereigniszuordnung oder
die Geometrie selbst quantisiert ist, kann die klassische Größe
\(\Delta s^2\) nicht ohne zusätzliche Brücke als von Anfang an scharfe
Zahl vorausgesetzt werden. Dann ist zunächst eine quantische oder
statistische Kausalstruktur zu modellieren, aus der der klassische
Lichtkegel hervorgeht.

\subsubsection{2.3 Planck-Skala}\label{planck-skala}

Die natürliche Planck-Normalform verbindet

\[
\frac{\ell_P}{t_P}=c,
\qquad
\frac{E_P}{p_P}=c,
\]

\[
\ell_Pp_P=\hbar,
\qquad
t_PE_P=\hbar,
\]

und am Planck-Punkt

\[
\ell_P=\frac{\hbar}{m_Pc}=\frac{Gm_P}{c^2}.
\]

Der letzte Ausdruck verwendet den Gravitationsradius \(Gm/c^2\), nicht
den um den Faktor zwei größeren Schwarzschild-Radius. Diese Identitäten
bestimmen eine exakte Maßstruktur. Sie liefern allein noch keine
mikroskopische Dynamik und keinen empirischen Nachweis eines diskreten
Raumzeitatoms.

\subsubsection{2.4 Schwarze Löcher und
Singularitäten}\label{schwarze-luxf6cher-und-singularituxe4ten}

In der Allgemeinen Relativitätstheorie werden Singularitäten fachlich
nicht primär als materielle Punkte definiert, sondern über das Scheitern
einer regulären Fortsetzung, insbesondere geodätische Unvollständigkeit
unter den jeweiligen Voraussetzungen. QCE ersetzt diesen klassischen
Grenzbefund nicht per Definition durch ein neues Naturgesetz. Es
formuliert eine Korrespondenzhypothese:

\[
\text{klassische Fortsetzungsgrenze}
\quad\rightsquigarrow\quad
\text{quantengravitative Übergangszone}.
\]

Modelle von Schwarz-zu-Weißloch-Übergängen, nichtsingulären Kernen und
quantisierten kausalen Strukturen zeigen, dass diese Frage ein
etablierter Forschungsgegenstand ist. Sie bestätigen jedoch nicht
automatisch die konkrete QCE-Zweischrittdynamik.

\begin{center}\rule{0.5\linewidth}{0.5pt}\end{center}

\subsection{3. Ontologie des
Unterschieds}\label{ontologie-des-unterschieds}

QCE übernimmt die QIK-VRT-Grundordnung

\[
D \rightarrow I \rightarrow R \rightarrow W \rightarrow C,
\]

wobei

\begin{itemize}
\tightlist
\item
  \(D\) einen typisierten Unterschied,
\item
  \(I\) die daraus gewinnbare Information,
\item
  \(R\) eine explizite Relation,
\item
  \(W\) eine beobachtbare oder erwartete Wirkung und
\item
  \(C\) eine nur unter angegebenen Brücken zulässige Kausalordnung
  bezeichnet.
\end{itemize}

Die Reihenfolge ist keine Behauptung, dass in der Natur stets fünf
diskrete Zeitpunkte in dieser Wortreihenfolge auftreten. Sie ist eine
Typordnung des Erklärens. Eine Kausalbehauptung darf erst aufgewertet
werden, wenn der relevante Unterschied, die Relation, die Wirkung und
die Evidenz explizit gebunden sind.

Der Satz

\begin{quote}
Kausalität ist Relation, nicht Sequenz.
\end{quote}

bestreitet nicht die Bedeutung zeitlicher Ordnung. Er bestreitet den
unzulässigen Schluss von bloßer zeitlicher Folge auf Ursache.

\begin{center}\rule{0.5\linewidth}{0.5pt}\end{center}

\subsection{4. Unsicherheitsbilanz statt „Entfernung`` der
Unschärfe}\label{unsicherheitsbilanz-statt-entfernung-der-unschuxe4rfe}

QCE zerlegt eine beobachtete Unsicherheitsbilanz schematisch in

\[
\Sigma_{\text{gesamt}}
=
\Sigma_{\text{instrumentell}}
+
\Sigma_{\text{grob}}
+
\Sigma_{\text{modell}}
+
\Sigma_{\text{irreduzibel}}.
\]

Die additive Schreibweise ist zunächst eine Modellvereinfachung. In
realen Auswertungen können Kreuzkovarianzen und nichtlineare
Abhängigkeiten auftreten. Der grundlegende methodische Anspruch bleibt:

\begin{enumerate}
\def\labelenumi{\arabic{enumi}.}
\tightlist
\item
  Instrumentelle Unsicherheit wird an Kalibration, Auflösung und
  Messprotokoll gebunden.
\item
  Statistische Unsicherheit wird an Stichprobe und Schätzverfahren
  gebunden.
\item
  Modellunsicherheit wird durch konkurrierende Modelle und
  Sensitivitätsanalysen ausgewiesen.
\item
  Grobkörnungsunsicherheit wird an die verlorene mikroskopische
  Auflösung gebunden.
\item
  Der verbleibende quantische Anteil wird nicht auf null gesetzt,
  sondern als irreduzibel erhalten.
\end{enumerate}

Der inverse Rekonstruktionsweg lautet

\[
D_N
\xrightarrow{\mathcal R_N^{-1}}
D_{N-1}
\xrightarrow{\mathcal R_{N-1}^{-1}}
\cdots
\xrightarrow{\mathcal R_2^{-1}}
D_1,
\]

wobei jede Inversion ihre Nicht-Eindeutigkeit, Informationsverluste und
Voraussetzungen dokumentieren muss. Eine nicht-injektive Grobkörnung
besitzt im Allgemeinen keine eindeutige Inverse. QCE verlangt deshalb
keinen magischen Rückgewinn verlorener Information. Zulässige Ergebnisse
sind:

\begin{itemize}
\tightlist
\item
  eindeutig rekonstruierbar,
\item
  äquivalenzklassenweise rekonstruierbar,
\item
  teilweise rekonstruierbar,
\item
  mehrdeutig,
\item
  evidenzbedingt offen,
\item
  irreversibel informationsverloren.
\end{itemize}

Damit wird „herausrechnen`` zu einer überprüfbaren Bilanz und nicht zu
einer Verletzung der Quantenmechanik.

\begin{center}\rule{0.5\linewidth}{0.5pt}\end{center}

\subsection{5. Der quantisch unscharfe
Kausalkegel}\label{der-quantisch-unscharfe-kausalkegel}

Auf fundamentaler Ebene verwendet QCE zunächst keine scharfe binäre
Kausalzuordnung. Schematisch wird ein Operator oder eine typisierte
Kausalvariable betrachtet,

\[
\widehat{s^2}
=
\widehat g_{\mu\nu}
\,\Delta\widehat x^\mu
\Delta\widehat x^\nu.
\]

Der Zustand liefert einen Erwartungswert und eine Streuung,

\[
\langle \widehat{s^2}\rangle,
\qquad
(\Delta s^2)^2
=
\langle (\widehat{s^2})^2\rangle
-
\langle \widehat{s^2}\rangle^2.
\]

Eine mögliche operationale Größe ist

\[
P_{\text{kausal}}
=
\operatorname{Tr}(\rho\,\Pi_{\text{zulässig}}),
\]

wobei \(\Pi_{\text{zulässig}}\) an den gewählten quantengravitativen
Formalismus gebunden werden muss. QCE behauptet nicht, dass diese
schematische Form bereits die eindeutige fundamentale Observable
darstellt. Sie markiert die Aufgabe: Vor dem klassischen Grenzfall wird
Kausalität als Zustands- und Relationsstruktur mit Unsicherheit
behandelt.

Der klassische Lichtkegel ist erreicht, wenn ein kontrollierter
Grenzübergang mindestens Folgendes liefert:

\begin{itemize}
\tightlist
\item
  eine effektiv klassische Metrik,
\item
  stabile Nullgrenzen,
\item
  hinreichend kleine relative Fluktuationen,
\item
  Mikrokausalität oder eine empirisch äquivalente Signalisierungsgrenze,
\item
  Lorentzverträglichkeit im erklärten Bereich,
\item
  robuste Übereinstimmung mit getesteten relativistischen Beobachtungen.
\end{itemize}

In Kurzform:

\[
\text{quantische Kausalrelation}
\xrightarrow{\text{Dekohärenz + Grobkörnung + Grenzwert}}
\text{klassischer Lichtkegel}.
\]

\begin{center}\rule{0.5\linewidth}{0.5pt}\end{center}

\subsection{6. Das
Planck-Übergangselement}\label{das-planck-uxfcbergangselement}

QCE bezeichnet mit \(\Sigma_P\) ein minimales Übergangselement des
Modells. Es ist ausdrücklich nicht als klassischer Punkt mit bereits
festgelegten Koordinaten definiert. Seine Bedeutung ist relational:

\[
\Sigma_P
=
\text{Träger eines elementaren Zustandsübergangs vor klassischer Geometrie}.
\]

Die physikalische Korrespondenzhypothese lautet:

\[
\text{Schwarze-Loch-Kernbereich}
\quad\longleftrightarrow\quad
\Sigma_P\text{-Dynamik}.
\]

Für diese Brücke sind mindestens erforderlich:

\begin{itemize}
\tightlist
\item
  eine kovariante Definition des Übergangselements,
\item
  eine Dynamik ohne vorausgesetzte klassische Zeitvariable oder eine
  erklärte relationale Uhr,
\item
  eine Energie-Impuls-Bilanz,
\item
  ein wohldefinierter Zustandsraum,
\item
  Stabilität und Unitarität im angegebenen Bereich,
\item
  eine Abbildung zu semiklassischen Schwarze-Loch-Observablen,
\item
  eine kontrollierte Behandlung des Horizonts,
\item
  eine Informationsbilanz über die gesamte Entwicklung.
\end{itemize}

Ohne diese Zeugen bleibt \(\Sigma_P\) ein formal definierter Kandidat.

\begin{center}\rule{0.5\linewidth}{0.5pt}\end{center}

\subsection{7. Sequenzielle
Zweischrittdynamik}\label{sequenzielle-zweischrittdynamik}

Der minimale QCE-Ablauf besteht aus zwei geordneten Übergängen:

\[
\mathcal H_{\Sigma,n}
\xrightarrow{U_{2n}}
\mathcal H_{\Sigma,n+\frac12}
\otimes \mathcal H_{a_n},
\]

\[
\mathcal H_{\Sigma,n+\frac12}
\otimes \mathcal H_{a_n}
\xrightarrow{U_{2n+1}}
\mathcal H_{\Sigma,n+1}
\otimes \mathcal H_{a_n}
\otimes \mathcal H_{b_n}.
\]

Der zusammengesetzte Übergang ist

\[
V_n=U_{2n+1}U_{2n}.
\]

Für einen informationserhaltenden Modellbereich wird mindestens eine
Isometrie verlangt,

\[
V_n^\dagger V_n=I,
\]

oder eine explizit begründete offene-System-Dynamik, deren
Gesamtentwicklung unitär ist.

Der erste Schritt materialisiert im Modell einen typisierten
Unterschied. Der zweite Schritt erzeugt nicht lediglich einen zweiten
unabhängigen Datensatz, sondern eine gemeinsame Paarrelation. Für einen
reinen Paarzustand kann eine Schmidt-Zerlegung geschrieben werden,

\[
|\Psi_n\rangle
=
\sum_k\sqrt{\lambda_k}
|k\rangle_{a_n}|k\rangle_{b_n},
\qquad
\sum_k\lambda_k=1.
\]

Mehr als ein von null verschiedener Schmidt-Koeffizient kennzeichnet
Verschränkung. Der beigefügte kleine Lean-Kern formalisiert nicht die
vollständige Hilbertraumtheorie dieser Gleichung. Er formalisiert die
endliche Modellgrenze: geordnete Schritte, gemeinsame Herkunft,
Paarbindung und die Tatsache, dass die physikalische
Verschränkungsbrücke ein eigener Zeuge bleibt.

\begin{center}\rule{0.5\linewidth}{0.5pt}\end{center}

\subsection{8. Globale Verschränkung statt isolierter
Paare}\label{globale-verschruxe4nkung-statt-isolierter-paare}

Ein Produkt unabhängiger Bell-Paare reicht nicht aus, um eine
zusammenhängende Raumzeit oder eine unitäre Schwarze-Loch-Evolution zu
begründen. Der neue Paarsektor muss mit der bereits vorhandenen
Relationsstruktur gekoppelt sein:

\[
\rho_{a_nb_nR_{<n}}
\neq
\rho_{a_nb_n}\otimes\rho_{R_{<n}}.
\]

Hier bezeichnet \(R_{<n}\) die vor dem Schritt vorhandenen relationalen
Freiheitsgrade. Die genaue Form der Nichtfaktorisierung, ihre
Monogamiebedingungen und ihre Dynamik müssen im vollständigen Modell
angegeben werden.

Für verdampfende Schwarze Löcher entsteht zusätzlich die
Informationsparadoxie. Eine unitäre Gesamtbeschreibung muss die
qualitative Page-Kurve reproduzieren: Die Entropie der Strahlung wächst
zunächst, erreicht ein Maximum und sinkt anschließend, wenn die
Gesamtstrahlung rein wird. QCE darf daher späte Quanten nicht
ausschließlich mit unabhängigen inneren Partnern verschränken. Die
globale Codierung muss die frühere Strahlung einbeziehen oder eine
äquivalente Informationsrekonstruktion liefern.

\texttt{PAGE\_CURVE\_CORRESPONDENCE} ist deshalb ein eigenständiges
Schließungsgate.

\begin{center}\rule{0.5\linewidth}{0.5pt}\end{center}

\subsection{9. Relationales Netz und
Raumzeitemergenz}\label{relationales-netz-und-raumzeitemergenz}

Nach \(N\) Zweischritten wird ein gewichtetes Relationsnetz betrachtet,

\[
G_N=(E_N,R_N,W_N),
\]

mit

\begin{itemize}
\tightlist
\item
  Ereignissen oder Zustandsmarkern \(E_N\),
\item
  kausalen und quantischen Relationen \(R_N\),
\item
  Zustands-, Phasen- oder Entropiegewichten \(W_N\).
\end{itemize}

Die Grobkörnung ist eine Abbildung

\[
\mathcal C:G_N\rightarrow(\mathcal M,g_{\mu\nu},\Phi),
\]

wobei \(\Phi\) effektive Materie- und Feldfreiheitsgrade bezeichnet.
Eine physikalische Raumzeitemergenz ist erst gezeigt, wenn
\(\mathcal C\) mindestens folgende Eigenschaften besitzt:

\begin{enumerate}
\def\labelenumi{\arabic{enumi}.}
\tightlist
\item
  \textbf{Dimensionskorrespondenz:} Der effektive Raum hat im relevanten
  Bereich die beobachtete Dimension und Signatur.
\item
  \textbf{Lorentzkorrespondenz:} Lokale Lorentzsymmetrie wird
  reproduziert oder Abweichungen werden quantitativ begrenzt.
\item
  \textbf{Einstein-Grenzfall:} Die effektive Dynamik liefert die
  getestete klassische Gravitation.
\item
  \textbf{Quantenfeld-Grenzfall:} Bekannte lokale Quantenfeldtheorie
  entsteht im geeigneten Bereich.
\item
  \textbf{Stabilität:} Keine unphysikalischen Geister oder
  unkontrollierten Instabilitäten treten im deklarierten Scope auf.
\item
  \textbf{Kausalität:} Signalisierungs- und Konsistenzbedingungen werden
  erfüllt.
\item
  \textbf{Nichtzirkularität:} Die Zielgeometrie wird nicht bereits in
  der mikroskopischen Definition versteckt vorausgesetzt.
\item
  \textbf{Empirische Unterscheidbarkeit:} Mindestens eine Beobachtung
  unterscheidet QCE von konkurrierenden Theorien.
\end{enumerate}

Die Forschung zu Verschränkung und Raumzeitkonnektivität, Entanglement
Equilibrium, kausalen Mengen und unbestimmter Kausalordnung liefert
wichtige Anschlussstellen. Keine dieser Arbeiten beweist allein die
konkrete QCE-Abbildung \(\mathcal C\).

\begin{center}\rule{0.5\linewidth}{0.5pt}\end{center}

\subsection{10. Bewusstsein als epistemische
Partition}\label{bewusstsein-als-epistemische-partition}

QCE trennt den physikalischen Zustandsraum von einer epistemischen
Klassifikation des Bewusstseins. Sei \(\Omega\) die Gesamtheit dessen,
was einem kognitiven System als Gegenstand der Verarbeitung zugänglich
oder denkbar wird. Dann wird eine disjunkte Partition verwendet:

\[
\Omega=W\,\dot\cup\,S\,\dot\cup\,R.
\]

Dabei bezeichnet

\begin{itemize}
\tightlist
\item
  \(W\): Wahrhaftigkeit - hinreichend stabilisierte, registrierte und
  evidenzgebundene Aussagen,
\item
  \(S\): Unsicherheit - explizit modellierte, noch nicht entschiedene
  Möglichkeiten,
\item
  \(R\): Rest - gegenwärtig nicht angemessen im Modell repräsentierte
  oder gänzlich unbekannte Möglichkeiten.
\end{itemize}

Diese Partition ist kein Ersatz für die Born-Regel und keine Behauptung,
Bewusstsein verursache den physikalischen Kollaps. Sie ist eine
Erkenntnis- und Verantwortungsordnung. Ihre Aufgabe ist, eine
Möglichkeit nicht als Tatsache und ein unbekanntes Außen nicht als bloß
bekannte Wahrscheinlichkeit auszugeben.

Der Satz

\begin{quote}
Bewusstsein ist Mengenlehre.
\end{quote}

ist innerhalb von QIK-VRT eine verdichtete interpretative These über die
Klassifikation von Erkenntniszuständen.

\begin{center}\rule{0.5\linewidth}{0.5pt}\end{center}

\subsection{11.
Round-Trip-Rekonstruktion}\label{round-trip-rekonstruktion}

QCE wird als inverses und anschließend vorwärts prüfbares Problem
formuliert:

\[
\text{Wirkung}
\rightarrow
\text{Evidenz}
\rightarrow
\text{Zustands- und Relationsrekonstruktion}
\rightarrow
\text{Modell}
\rightarrow
\text{formale Folgerung}
\rightarrow
\text{Vorhersage}
\rightarrow
\text{erneute Messung}.
\]

Der Round Trip gilt nur dann als geschlossen, wenn

\begin{enumerate}
\def\labelenumi{\arabic{enumi}.}
\tightlist
\item
  die Beobachtungen quellen- und messgebunden sind,
\item
  die inverse Rekonstruktion ihre Mehrdeutigkeit ausweist,
\item
  das Modell exakt versioniert ist,
\item
  die formalen Folgerungen reproduzierbar geprüft sind,
\item
  eine neue oder zumindest unabhängige Vorhersage abgeleitet wird,
\item
  die Zuordnung zu Messdaten vor der Auswertung hinreichend eingefroren
  ist,
\item
  die Naturbeobachtung mit Unsicherheitsbudget erfolgt,
\item
  unabhängige Reproduktion möglich ist.
\end{enumerate}

Ein formaler Kreis innerhalb eines selbstdefinierten Modells ist noch
kein physikalischer Round Trip.

\begin{center}\rule{0.5\linewidth}{0.5pt}\end{center}

\subsection{12. Der endliche
Lean-Modellvertrag}\label{der-endliche-lean-modellvertrag}

Der beigefügte Lean-Kandidat verwendet bewusst einen kleinen,
\texttt{Std}-basierten Kern. Er formalisiert:

\begin{itemize}
\tightlist
\item
  ein typisiertes Planck-Übergangselement,
\item
  eine kanonische Zweischrittspur,
\item
  Erhaltung der gemeinsamen Quellidentität,
\item
  eine endliche Paarrelationsdarstellung,
\item
  eine explizite Unsicherheitsbilanz,
\item
  Erhaltung der irreduziblen Komponente nach Entfernung modellierter
  reduzierbarer Anteile,
\item
  einen quantisch unaufgelösten Kausalstatus,
\item
  ein dreiteiliges Gate für einen klassischen Lichtkegel,
\item
  monotone Erweiterung eines endlichen Relationsnetzes,
\item
  eine konjunktive physikalische Schließungsregel,
\item
  eine Trennung von Kernel-Receipt und empirischer Korrespondenz.
\end{itemize}

Der Lean-Kern beweist nicht:

\begin{itemize}
\tightlist
\item
  dass reale Singularitäten Planck-Elemente sind,
\item
  dass reale Schwarze Löcher die QCE-Zweischrittdynamik ausführen,
\item
  dass die modellierte Paarrelation bereits physikalische Verschränkung
  im vollständigen Hilbertraumsinn ist,
\item
  dass die Grobkörnung die Einstein-Gleichungen ergibt,
\item
  dass die Page-Kurve reproduziert wird,
\item
  dass eine neue QCE-Vorhersage bestätigt wurde,
\item
  dass die historische Prioritätsbewertung feststeht.
\end{itemize}

Diese Grenzen werden in \texttt{CLAIM\_MATRIX.json},
\texttt{SOURCE\_EVIDENCE\_BINDINGS.json} und im Kernel-Receipt
festgehalten.

\begin{center}\rule{0.5\linewidth}{0.5pt}\end{center}

\subsection{13. Fail-closed Schließung}\label{fail-closed-schlieuxdfung}

Der QCE-Kandidat gilt nur dann als physikalisch geschlossen, wenn
sämtliche folgenden Zeugen vorliegen:

\begin{enumerate}
\def\labelenumi{\arabic{enumi}.}
\tightlist
\item
  Planck-Skala-Korrespondenz,
\item
  explizite Zweischrittdynamik,
\item
  physikalische Paarverschränkung,
\item
  globale Verschränkungsintegration,
\item
  vollständige Unsicherheitsbilanz,
\item
  Unitarität oder äquivalente Gesamtinformationserhaltung,
\item
  Energie-Impuls-Erhaltung,
\item
  Page-Kurven-Korrespondenz für den relevanten Schwarze-Loch-Scope,
\item
  Quantenfeld-Grenzfall,
\item
  klassischer Einstein-Grenzfall,
\item
  klassischer Lichtkegel-Grenzfall,
\item
  kausale Konsistenz,
\item
  Nichtzirkularität,
\item
  falsifizierbare unterscheidende Vorhersage,
\item
  empirische Korrespondenz,
\item
  unabhängige Reproduktion.
\end{enumerate}

Formal:

\[
\mathrm{QCEClosed}
=
\bigwedge_{i=1}^{16} W_i.
\]

Fehlt ein Zeuge, bleibt der Gesamtstatus offen. Ein Kernel-Exitcode null
ersetzt keinen der empirischen oder physikalischen Zeugen.

\begin{center}\rule{0.5\linewidth}{0.5pt}\end{center}

\subsection{14. Falsifikationsprogramm}\label{falsifikationsprogramm}

Ein wissenschaftlicher QCE-Kandidat benötigt quantitative Abweichungen
oder neue Korrelationen. Mögliche Arbeitsrichtungen sind:

\subsubsection{14.1 Ringdown- und
Echo-Struktur}\label{ringdown--und-echo-struktur}

Ein nichtklassischer Kern oder eine Übergangszone könnte,
modellabhängig, Abweichungen in späten Ringdown-Signalen erzeugen. Ein
belastbarer Test verlangt eine vorab spezifizierte Wellenformfamilie,
systematische Fehler, Detektorrauschen und einen Vergleich mit
konventionellen Umgebungs- und Modellunsicherheiten.

\subsubsection{14.2 Hawking- und
Page-Kurven-Korrespondenz}\label{hawking--und-page-kurven-korrespondenz}

Die Zweischrittdynamik muss eine Informationsstruktur liefern, die mit
unitärer Verdampfung vereinbar ist. Eine reine Paarerzeugung ohne
globale Rekodierung wäre ausgeschlossen.

\subsubsection{14.3
Kausalkegel-Fluktuationen}\label{kausalkegel-fluktuationen}

Falls QCE residuale Kausalstrukturfluktuationen vorhersagt, müssen
Größe, Spektrum, Skalierung und Signalisierungsgrenze quantitativ
berechnet werden. Bestehende experimentelle Schranken dürfen nicht
verletzt werden.

\subsubsection{14.4 Lorentz- und
Dispersionsgrenzen}\label{lorentz--und-dispersionsgrenzen}

Eine diskrete oder relationale Mikrodynamik darf keine bereits
ausgeschlossenen energieabhängigen Laufzeit- oder Dispersionssignaturen
erzeugen. Entweder wird lokale Lorentzinvarianz emergent exakt
reproduziert oder verbleibende Abweichungen werden unter bestehende
Schranken gedrückt.

\subsubsection{14.5
Entropie-Flächen-Beziehung}\label{entropie-fluxe4chen-beziehung}

Das Relationsnetz sollte eine nachvollziehbare Verbindung zwischen
Verschränkungsentropie und effektiver Geometrie liefern. Ein bloßes
Einsetzen einer bekannten Flächenformel wäre zirkulär.

Keine dieser Richtungen ist im vorliegenden endlichen Lean-Kern
geschlossen.

\begin{center}\rule{0.5\linewidth}{0.5pt}\end{center}

\subsection{15. Gegenmodelle und mögliche
Fehlerklassen}\label{gegenmodelle-und-muxf6gliche-fehlerklassen}

QCE muss insbesondere gegen folgende Fehlerklassen geprüft werden:

\begin{itemize}
\tightlist
\item
  \textbf{Begriffsäquivokation:} „Singularität``, „Planck-Element``,
  „Paar``, „Verschränkung`` oder „Kollaps`` werden zwischen Modell und
  Natur unbemerkt umgedeutet.
\item
  \textbf{Zirkularität:} Raumzeit oder Lichtkegel werden in der
  Mikrodynamik vorausgesetzt und anschließend als emergent ausgegeben.
\item
  \textbf{Dimensionsanalyse ohne Dynamik:} Exakte Planck-Identitäten
  werden fälschlich als vollständige Feldgleichungen behandelt.
\item
  \textbf{Unschärfevernichtung:} Eine Rekonstruktion wird unzulässig als
  Beseitigung irreduzibler Quantenfluktuation ausgegeben.
\item
  \textbf{Lokale Paarfalle:} Unabhängige Paare werden ohne globale
  Informationsstruktur als Lösung des Informationsparadoxons behauptet.
\item
  \textbf{Kernel-Promotion:} Formale Typ- und Gate-Theoreme werden als
  empirischer Naturbeweis ausgegeben.
\item
  \textbf{Unterbestimmtheit:} Mehrere Mikromodelle erzeugen dieselben
  makroskopischen Daten, ohne dass QCE eine unterscheidende Vorhersage
  liefert.
\item
  \textbf{Nichtreproduzierbarkeit:} Quellen, Parameter, Code, Toolchain
  oder Ausführung sind nicht bytegenau gebunden.
\item
  \textbf{Prioritätsüberdehnung:} Eine persönliche Größenbewertung wird
  als fachlicher Konsens dargestellt.
\end{itemize}

Die QIK-VRT-Prüfarchitektur soll diese Fehler nicht rhetorisch, sondern
maschinenlesbar blockieren.

\begin{center}\rule{0.5\linewidth}{0.5pt}\end{center}

\subsection{16. Repository-, Lean- und
Zenodo-Workflow}\label{repository--lean--und-zenodo-workflow}

Der vorgesehene Ablauf ist serialisiert:

\begin{enumerate}
\def\labelenumi{\arabic{enumi}.}
\tightlist
\item
  Aktuelle Authority- und Mirror-Mains werden neu beobachtet.
\item
  Der aktuell autorisierte Repository-Arbeitsauftrag wird abgeschlossen
  oder ausdrücklich neu priorisiert.
\item
  Der QCE-Kandidat wird additiv auf einem neuen Authority-Branch
  materialisiert.
\item
  Der read-only QIK-VRT-Bootloader wird ausgeführt.
\item
  Python-Validator, Negativtests, Dokumentprüfung und Prüfsummen laufen.
\item
  Lean 4.19.0 kompiliert den Modellkern und führt den Axiom-Audit aus.
\item
  Ein source-, commit-, tree-, toolchain- und output-gebundenes
  Kernel-Receipt wird erzeugt.
\item
  Exakte Head-Gates und Repository-Integrität laufen.
\item
  Erst danach wird ein byteidentischer Mirror-Kandidat erzeugt und
  separat geprüft.
\item
  Nach Review und Promotion wird ein Zenodo-Kandidat aus den
  freigegebenen Bytes erzeugt.
\item
  Die Zenodo-Veröffentlichung verlangt eine ausdrückliche
  Owner-Autorisierung und eine erneute Hash-Prüfung.
\item
  Nach Veröffentlichung wird ein Rücklauf-Receipt mit DOI, Record-ID,
  Dateihashes und Repository-Bindung persistiert.
\end{enumerate}

GitHub ist die versionierte Ausführungs- und Reviewebene. Lean ist die
formale Prüfebene. Zenodo ist die persistente Publikations- und
Zitierungsebene. Keine Ebene ersetzt die andere.

\begin{center}\rule{0.5\linewidth}{0.5pt}\end{center}

\subsection{17. Gegenwärtiger Status}\label{gegenwuxe4rtiger-status}

Der vorliegende Lieferstand ist ein \textbf{commit-ready
Kandidatenpaket}. Markdown, wissenschaftliche Spezifikation,
Lean-Quellkandidat, Validatoren, Workflow, Claim-Matrix,
Quellenbindungen, Zenodo-Metadaten und Prüfsummen werden gemeinsam
bereitgestellt.

Der neue Lean-Quellkandidat wurde in dieser Ausführungsumgebung noch
nicht mit Lean 4.19.0 kernelgeprüft. Deshalb enthält das Paket
ausschließlich ein \texttt{KERNEL\_RECEIPT\_TEMPLATE.json} mit dem
Status \texttt{NOT\_EXECUTED}. Ein späterer erfolgreicher
GitHub-Actions-Lauf muss dieses Template durch ein exakt gebundenes
Ausführungsreceipt ersetzen.

Die physikalische Korrespondenz bleibt unabhängig davon offen.

\texttt{FORMAL\_SOURCE\ =\ PREPARED}

\texttt{LEAN\_EXECUTION\ =\ PENDING\_REPOSITORY\_RUN}

\texttt{PHYSICAL\_CORRESPONDENCE\ =\ OPEN\_CANDIDATE}

\texttt{ZENODO\_PUBLICATION\ =\ NOT\_EXECUTED}

\texttt{PASS\ =\ NOT\_CLAIMED}

\texttt{FINAL\_PASS\ =\ NOT\_CLAIMED}

\texttt{EFFECT\_ACK\_DONE\ =\ NOT\_CLAIMED}

\texttt{EFFECT\_STATE\ =\ EFFECT\_ACK\_CONTINUE}

\begin{center}\rule{0.5\linewidth}{0.5pt}\end{center}

\subsection{18. Schlussfolgerung}\label{schlussfolgerung}

QCE ordnet eine weitreichende Hypothese in einen prüfbaren
Forschungsvertrag. Der klassische Lichtkegel wird nicht verworfen,
sondern als makroskopische Grenzstruktur einer tieferen quantischen
Kausalordnung interpretiert. Die Heisenbergsche Unschärfe wird nicht
entfernt, sondern durch die Rekonstruktion bilanziert und bis zum
irreduziblen Rest erhalten. Eine klassische Singularität wird als
mögliche Projektionsgrenze einer quantengravitativen Übergangszone
behandelt. Ein sequenzieller Zweischritt erzeugt im Modell zunächst
einen Unterschied und anschließend eine relationale Paarstruktur. Viele
global gebundene Relationen bilden den Kandidaten für eine emergente
Raumzeit.

Die konzeptionelle Verdichtung ist stark. Ihre physikalische Gültigkeit
hängt jedoch an expliziten Brücken, die nicht durch Benennung ersetzt
werden dürfen. Der beigefügte formale Kern soll genau diese Grenze
durchsetzen: Er macht die interne Modelllogik maschinenprüfbar und
blockiert zugleich die unzulässige Promotion eines Kernelbeweises zur
experimentellen Entdeckung.

Der wissenschaftliche Anspruch ist damit weder bloße Spekulation noch
bereits abgeschlossene Weltformel. Er ist ein reproduzierbares Programm:

\[
\text{Unterschied}
\rightarrow
\text{Relation}
\rightarrow
\text{quantische Kausalstruktur}
\rightarrow
\text{klassische Raumzeit}
\rightarrow
\text{Vorhersage}
\rightarrow
\text{Messung}.
\]

Der Round Trip entscheidet.

\begin{center}\rule{0.5\linewidth}{0.5pt}\end{center}

\subsection{Literatur}\label{literatur}

\begin{enumerate}
\def\labelenumi{\arabic{enumi}.}
\tightlist
\item
  W. Heisenberg, „Über den anschaulichen Inhalt der quantentheoretischen
  Kinematik und Mechanik``, \emph{Zeitschrift für Physik} 43 (1927),
  172-198, DOI: 10.1007/BF01397280.
\item
  H. P. Robertson, „The Uncertainty Principle``, \emph{Physical Review}
  34 (1929), 163-164, DOI: 10.1103/PhysRev.34.163.
\item
  S. W. Hawking und R. Penrose, „The Singularities of Gravitational
  Collapse and Cosmology``, \emph{Proceedings of the Royal Society A}
  314 (1970), 529-548, DOI: 10.1098/rspa.1970.0021.
\item
  S. W. Hawking, „Particle Creation by Black Holes``,
  \emph{Communications in Mathematical Physics} 43 (1975), 199-220, DOI:
  10.1007/BF02345020.
\item
  S. Doplicher, K. Fredenhagen und J. E. Roberts, „The Quantum Structure
  of Spacetime at the Planck Scale and Quantum Fields``,
  \emph{Communications in Mathematical Physics} 172 (1995), 187-220,
  DOI: 10.1007/BF02104515.
\item
  D. Malament, „The Class of Continuous Timelike Curves Determines the
  Topology of Spacetime``, \emph{Journal of Mathematical Physics} 18
  (1977), 1399-1404, DOI: 10.1063/1.523436.
\item
  L. Bombelli, J. Lee, D. Meyer und R. D. Sorkin, „Space-Time as a
  Causal Set``, \emph{Physical Review Letters} 59 (1987), 521-524, DOI:
  10.1103/PhysRevLett.59.521.
\item
  C. Oreshkov, F. Costa und Č. Brukner, „Quantum Correlations with No
  Causal Order``, \emph{Nature Communications} 3 (2012), 1092, DOI:
  10.1038/ncomms2076.
\item
  M. Van Raamsdonk, „Building up Spacetime with Quantum Entanglement``,
  \emph{General Relativity and Gravitation} 42 (2010), 2323-2329, DOI:
  10.1007/s10714-010-1034-0.
\item
  T. Jacobson, „Entanglement Equilibrium and the Einstein Equation``,
  \emph{Physical Review Letters} 116 (2016), 201101, DOI:
  10.1103/PhysRevLett.116.201101.
\item
  J. Maldacena und L. Susskind, „Cool Horizons for Entangled Black
  Holes``, \emph{Fortschritte der Physik} 61 (2013), 781-811, DOI:
  10.1002/prop.201300020.
\item
  A. Almheiri, D. Marolf, J. Polchinski und J. Sully, „Black Holes:
  Complementarity or Firewalls?{}``, \emph{Journal of High Energy
  Physics} 2013, 62, DOI: 10.1007/JHEP02(2013)062.
\item
  G. Penington, „Entanglement Wedge Reconstruction and the Information
  Paradox``, \emph{Journal of High Energy Physics} 2020, 2, DOI:
  10.1007/JHEP09(2020)002.
\item
  C. Rovelli und F. Vidotto, „Planck Stars``, \emph{International
  Journal of Modern Physics D} 23 (2014), 1442026, DOI:
  10.1142/S0218271814420267.
\item
  J. F. Donoghue, „General Relativity as an Effective Field Theory: The
  Leading Quantum Corrections``, \emph{Physical Review D} 50 (1994),
  3874-3888, DOI: 10.1103/PhysRevD.50.3874.
\item
  Lean FRO, \emph{Lean Language Reference} und \emph{Validating a Lean
  Proof}, jeweils für die im Repository gebundene Toolchain und
  Kernelgrenze.
\end{enumerate}

\begin{center}\rule{0.5\linewidth}{0.5pt}\end{center}

\textbf{Quod erat demonstrandum gilt nur für die ausdrücklich benannten
und tatsächlich ausgeführten formalen Modelltheoreme.}

\textbf{Ingolf Lohmann}

\end{document}
